% =====================================================================
%  Thème 4 — Équations trigonométriques — NIVEAU MOYEN
%  Exercices
% =====================================================================
\documentclass[11pt,a4paper]{article}
\usepackage[utf8]{inputenc}
\usepackage[T1]{fontenc}
\usepackage{lmodern}
\usepackage[french]{babel}
\usepackage{amsmath,amssymb,mathtools}
\usepackage{icomma}
\usepackage{xcolor}
\usepackage{colortbl}
\usepackage{tikz}
\usepackage[most]{tcolorbox}
\usepackage{enumitem}
\usepackage{array}
\usepackage[a4paper,margin=2.2cm]{geometry}
\usetikzlibrary{arrows.meta,calc}

\definecolor{primary}{RGB}{214,51,108}
\definecolor{primarydark}{RGB}{140,20,64}

\DeclareMathOperator{\tg}{tg}
\DeclareMathOperator{\cotg}{cotg}
\newcommand{\degr}{^{\circ}}
\newcommand{\Z}{\mathbb{Z}}

\newtcolorbox{consigne}{enhanced,breakable,colback=primary!5,colframe=primary,
  boxrule=0.8pt,arc=2pt,left=3mm,right=3mm,top=2mm,bottom=2mm}

\newcounter{exo}
\newcommand{\exo}[1]{\medskip\refstepcounter{exo}%
  \noindent{\large\bfseries\color{primarydark}Exercice \theexo}%
  \ \textcolor{primary}{\textbullet}\ \textbf{#1}\par\nopagebreak\smallskip}

\setlength{\parindent}{0pt}
\pagestyle{empty}

\begin{document}

\begin{center}
{\LARGE\bfseries\color{primary}Thème 4 — Équations trigonométriques}\\[2pt]
{\color{primarydark}\textsc{Niveau Moyen — Exercices}}
\end{center}
\hrule height 1pt
\medskip

\begin{consigne}
\textbf{Objectif :} gérer les \textbf{valeurs négatives} (choix du bon angle de référence et du bon
quadrant), isoler le nombre trigonométrique dans une équation du type $a\cdot\sin x+b=0$, et
\textbf{sélectionner} les solutions imposées par un intervalle ou un quadrant. \textbf{Sans
calculatrice}, valeurs exactes, conclure par $S$.
\end{consigne}

\exo{Valeurs négatives}
Résous dans $\mathbb{R}$ :
\[
\textbf{(a)}\ \cos x=-\frac12 \qquad\qquad
\textbf{(b)}\ \sin x=-\frac{\sqrt3}{2}.
\]

\exo{Isoler puis sélectionner (cosinus)}
On considère l'équation $2\cos x+1=0$, avec $x$ situé dans le \textbf{2\textsuperscript{e} quadrant}.
Résous-la et donne la (les) solution(s) retenue(s).

\exo{Isoler puis sélectionner (sinus)}
On considère l'équation $2\sin x-\sqrt3=0$, avec $x$ situé dans le \textbf{2\textsuperscript{e}
quadrant}. Résous-la et donne la (les) solution(s) retenue(s).

\exo{Sélection sur un intervalle (tangente)}
Résous dans $[0\,;\,2\pi[$ :
\[
\tg x=-1 .
\]

\exo{Sélection sur un intervalle (sinus)}
Résous dans $[0\,;\,2\pi[$ :
\[
\sin x=-\frac12 .
\]

\end{document}
